\section{Six obstructions, with counterexamples}
\label{sec:breaks}

This section states the obstructions that any argument from the
fixed-spacing construction of \S\ref{sec:wilson} to a continuum mass gap
must clear. They were produced by a deliberate adversarial audit of this
program's own drafts, and four of them killed arguments that had been
written down here as complete.

They are stated as mathematics, not as caveats. Five carry explicit
finite counterexamples, and a finite counterexample to a proof strategy
is a permanent fact about that strategy. A reader who wants to know
where the remaining difficulty of the problem actually sits should read
this section and \S\ref{sec:obligations} and can safely skip everything
else in Part~\ref{part:continuum}.

The counterexamples are not rhetorical. Those of
Break~\ref{sec:break4} are registered checks in this repository's own
test suite --- \srcpath{tests/test\_os\_kernel\_gap.py}, six checks,
passing --- and the two collapse conditions of Break~\ref{sec:break3}
are checked in \srcpath{tests/test\_moving\_time\_gap.py}. An
obstruction that a repository asserts against itself, and then runs in
continuous integration, is harder to forget than one recorded in prose.

\subsection{Break 1: an \texorpdfstring{$O(g_k^2)$}{O(g\_k\textasciicircum 2)} increment bound does not give convergence}
\label{sec:break1}

A renormalization-group argument for the continuum limit produces a
sequence of effective actions $S_k^{\mathrm{eff}}$ and observable
expectations $m_k(\mathcal{O})$ indexed by scale, and controls the
one-step increment by the running coupling:
\begin{equation}
  \abs{m_{k+1}(\mathcal{O}) - m_k(\mathcal{O})}
  \;\le\; C_{\mathcal{O}}\, g_k^2
  \;\le\; \frac{C_{\mathcal{O}}'}{k + k_0},
  \label{eq:break1-bound}
\end{equation}
the second inequality being asymptotic freedom, $g_k^2 = \kappa/(k+k_0)$.
The step that must not be taken is to conclude from
\eqref{eq:break1-bound} that $(m_k)$ is Cauchy. Telescoping gives only
\begin{equation}
  \abs{m_n(\mathcal{O}) - m_m(\mathcal{O})}
  \;\le\; C_{\mathcal{O}}' \sum_{k=m}^{n-1} \frac{1}{k+k_0}
  \;\asymp\; C_{\mathcal{O}}' \log\!\left(\frac{n}{m}\right),
\end{equation}
which does not tend to zero as $m, n \to \infty$.

\begin{proposition}[Counterexample]
\label{prop:break1}
Let $x_k = \sin\!\big(\log(k+1)\big)$. Then $x_k \in [-1,1]$ and, by the
mean value theorem, $\abs{x_{k+1} - x_k} \le 1/(k+1) \to 0$, yet the set
of subsequential limits of $(x_k)$ is the whole interval $[-1,1]$.
\end{proposition}

So the increment bound \eqref{eq:break1-bound} is consistent with the
observable expectations oscillating forever across a fixed interval. What
is needed is one more power: $\abs{m_{k+1} - m_k} = O(g_k^4) = O(k^{-2})$.

The obstruction to obtaining it is not technical. On the asymptotically
free trajectory $u_k = (\alpha + \beta k)^{-1}$,
\begin{equation}
  \Delta u_k = -\frac{\beta}{(\alpha+\beta k)\big(\alpha + \beta(k+1)\big)},
  \qquad
  \frac{\abs{\Delta u_k}}{4 u_k u_{k+1}} = \frac{\beta}{4} = O(1),
\end{equation}
and differentiating an expectation along the trajectory,
\begin{equation}
  \frac{d}{ds}\,\mu_s(\mathcal{O}_s)
  = \mu_s(\partial_s \mathcal{O}_s)
  - \operatorname{Cov}_{\mu_s}\!\big(\mathcal{O}_s, \partial_s S_s\big),
  \label{eq:break1-cov}
\end{equation}
the normalization carries no small factor. The natural size of the
variation is therefore $O(g_k^2)$, and a bounded-covariance estimate
cannot improve it. Converting $1/k$ into $1/k^2$ requires an exact
cancellation of the leading partition-function covariance in
\eqref{eq:break1-cov}. No such cancellation is proved here.

\begin{scopenote}
This obstruction applies to the multiscale route recorded as\\
\texttt{\small DERIV:YANGMILLS\_CONTINUUM\_\allowbreak{}BALABAN\_\allowbreak{}MULTISCALE\_\allowbreak{}PROOF:\allowbreak{}THEOREM\_7\_1},\\
whose Cauchy input is exactly what \eqref{eq:break1-bound} fails to
supply. It does not bear on any statement in
Part~\ref{part:exact} or \ref{part:fixed}, which are at fixed scale.
\end{scopenote}

\subsection{Break 2: the polymer activity and the mass schedule are algebraically incompatible}
\label{sec:break2}

A multiscale polymer cluster expansion carries an activity parameter
$\kappa_k$ that must contract for the expansion to converge, say
\begin{equation}
  \kappa_{k+1} = \tfrac{3}{4}\,\kappa_k,
  \qquad \kappa_k = \kappa_0 \left(\tfrac{3}{4}\right)^{k} \longrightarrow 0,
\end{equation}
while a physical mass is extracted at each scale by
\begin{equation}
  \frac{\kappa_k - \log C_4}{a_k} = c_0 \Lambda > 0
  \quad\text{for all } k \ge 0,
  \qquad a_k = a_0 L^{-k},\ L = 3,
  \label{eq:break2-mass}
\end{equation}
with $C_4 > 1$ the polymer coordination constant.

\begin{proposition}
\label{prop:break2}
The two displays above are incompatible. Since $\kappa_k \to 0$ and
$\log C_4 > 0$, there is a finite
\[
  k_* = \left\lceil \frac{\log\!\big(\kappa_0/\log C_4\big)}{\log(4/3)} \right\rceil
\]
with $\kappa_k - \log C_4 < 0$ for every $k > k_*$; beyond that scale the
cluster sum $\sum_n e^{-(\kappa_k - \log C_4)n}$ diverges and the
expansion does not exist. Setting $C_4 = 1$, which discards polymer
entropy, does not repair it: then
$\kappa_k / a_k = (\kappa_0/a_0)(3L/4)^k = (\kappa_0/a_0)(9/4)^k \to \infty$,
so \eqref{eq:break2-mass} cannot equal a finite constant either.
\end{proposition}

The content of Proposition~\ref{prop:break2} is that a contracting
activity and a scale-independent physical mass cannot both be read off
the same schedule. Any repair must decouple them --- the mass must come
from a quantity that is not the activity divided by the spacing.

\subsection{Break 3: the moving-time gap can be exactly zero}
\label{sec:break3}

The moving-time theorem of \S\ref{sec:continuum} extracts a limiting
gap from approximate sources and one growing observation time per
cutoff. Under the error bound
$K_n(t) \le A a_n^{-p} e^{-mt} + B a_n^{r}$ and probe tilt
$\abs{\alpha_n} \ge c_\alpha a_n^{s}$, it gives
\begin{equation}
  M_* = \frac{m\,(r - 2s)}{p + r}
  \qquad\text{at}\qquad
  t_n = \frac{p+r}{m}\,\log\frac{1}{a_n}.
  \label{eq:break3}
\end{equation}
Two ways this returns nothing.

\emph{Zero margin.} If $r \le 2s$ then $M_* \le 0$: the theorem is
sharp and it yields exactly zero. Here $r$ is the convergence rate of the
coarse-graining background and $s$ the rate at which probes must shrink
to stay in the quadratic basin of the Wilson action. No calculation in
this program, or to this author's knowledge anywhere, certifies
$r > 2s$ for four-dimensional Yang--Mills. The inequality $r > 2s$ is
therefore a hypothesis with no current route, not a technical condition.

\emph{Finite horizon.} If the observation time is capped by the box,
$t_n \le c_{\max}\log(1/a_n)$, the optimum becomes
\begin{equation}
  M_{\mathrm{opt}} = \min\!\left(m - \frac{p+2s}{c_{\mathrm{opt}}},\
  \frac{r-2s}{c_{\mathrm{opt}}}\right),
  \qquad
  c_{\mathrm{opt}} = \min\!\left(c_{\max},\ \frac{p+r}{m}\right),
\end{equation}
and $c_{\max} \le (p+2s)/m$ forces $M_{\mathrm{opt}} \le 0$. Both
collapses are registered checks rather than remarks.

\subsection{Break 4: a local probe can be blind to the lightest state}
\label{sec:break4}

Suppose one measures a plaquette correlator and observes
\begin{equation}
  \ip{\mathcal{O}_{\mathrm{plaq}}}{e^{-tH}\mathcal{O}_{\mathrm{plaq}}}
  \le C e^{-M_{\mathrm{obs}} t},
  \qquad M_{\mathrm{obs}} = \log 4 \approx 1.386 .
\end{equation}
This does not bound the spectrum.

\begin{proposition}[Dark sector]
\label{prop:break4}
On $\C^3$ let
\[
  T = \operatorname{diag}\!\left(1,\ \tfrac{15}{16},\ \tfrac14\right),
  \qquad \Omega = e_1,\qquad v_{\mathrm{obs}} = e_3 .
\]
Then $\ip{v_{\mathrm{obs}}}{T^t v_{\mathrm{obs}}} = 4^{-t} = e^{-t\log 4}$,
while the true first excited energy is
$-\log(15/16) \approx 0.0645$, smaller by a factor above $21$. The probe
has zero overlap with the state realizing the gap.
\end{proposition}

Consequently a decay rate for local plaquettes is a statement about the
gap only after the plaquette family is proved \emph{total} in
$\Omega^{\perp}$. This is why totality is a conclusion, not a hypothesis,
in the fixed-spacing band theorem (\S\ref{sec:wilson}); and it is why a
continuum argument may not inherit totality from the fixed-spacing
statement without proving it again in the limit.

A second, subtler version of the same failure concerns positivity rather
than overlap.

\begin{proposition}[Osterwalder--Schrader null states]
\label{prop:break4b}
Let $(x,y) \in \{-1,1\}^2$ be uniform and let time reflection act by
$\Theta(x,y) = (y,x)$. The centered observable $F(x,y) = y$ has Euclidean
variance $\mathbb{E}[F^2] = 1 > 0$ but physical norm
$\norm{F}_{\mathrm{OS}}^2 = \ip{\Theta F}{F} = \mathbb{E}[xy] = 0$.
\end{proposition}

So positive Euclidean variance does not certify that a probe exists in
the physical Hilbert space at all. An algorithm that selects probes by
Euclidean variance can select states that vanish after reconstruction.

A third version concerns the generator rather than the states: with
$H_n = \operatorname{diag}(0,n)$ one has
$e^{-tH_n} \to \operatorname{diag}(1,0)$ for every $t>0$, so the limiting
semigroup is not strongly continuous at $t=0$ and no self-adjoint
generator exists on the second coordinate. Convergence of Gram matrices
alone therefore does not produce a limiting Hamiltonian.

\subsection{Break 5: the proved regime and the continuum regime are disjoint}
\label{sec:break5}

This is the structural obstruction, and the one that most changes what a
correct argument must look like.

The fixed-spacing band theorem (\S\ref{sec:wilson}) holds on
\begin{equation}
  \abs{u} \;\le\; \frac{u_*}{10\,022\,400\,000\,N},
  \label{eq:break5-u}
\end{equation}
and the spatial-closure estimates hold for
$0 \le \lambda \le \lambda_c$ with
$\lambda_c \in \left(\tfrac{1}{73}, \tfrac{1}{72}\right)$. Both are
neighbourhoods of zero in the strong-coupling variable.

The continuum limit lives elsewhere. With
$g^{-2}(a) = \beta_0 \log(1/a\Lambda)$ and $u = 2/\big(g^2(a)N\big)$,
\begin{equation}
  u(a) \longrightarrow \infty,
  \qquad
  \lambda(a) \sim g^{-4}(a) \longrightarrow \infty
  \qquad (a \to 0).
\end{equation}
So the interval on which the theorem is proved and the trajectory along
which the continuum is taken have empty intersection:
\begin{equation}
  \left[0, \lambda_c\right] \,\cap\, \left[\lambda(a), \infty\right) = \emptyset
  \quad\text{for all sufficiently small } a .
\end{equation}

\begin{scopenote}
The fixed-spacing results establish a strictly positive gap at fixed
lattice spacing in a neighbourhood of strong coupling. They supply no
control on the weak-coupling trajectory. An inference from
\eqref{eq:break5-u} to a continuum mass gap is a regime error, and this
program has made and retracted a version of it.
\end{scopenote}

What Break~5 demands of a future argument is a \emph{controlled
trajectory}: an estimate that survives continuation in the coupling, not
a better constant at $u \approx 0$. Improving the constant
$10\,022\,400\,000$ in \eqref{eq:break5-u} by any finite factor does not
move the boundary of the difficulty. This is why Problem~\ref{prob:trajectory}
in \S\ref{sec:obligations}, and not a sharper fixed-spacing estimate, is
the load-bearing open problem.

\subsection{Break 6: diffusion time is not Euclidean time}
\label{sec:break6}

A recurring route derives a spectral gap from curvature: positive
Bakry--\'Emery--Ricci curvature on the configuration manifold gives a
logarithmic Sobolev inequality, which gives a spectral gap
$\lambda_{\mathrm{diff}} > 0$ for the Langevin generator. The step that
fails is the identification of $\lambda_{\mathrm{diff}}$ with a physical
mass.

Langevin time is a fictitious parabolic parameter; Euclidean time is a
spacetime coordinate of the reconstructed theory. In free models the two
are related by $\Delta_{\mathrm{phys}} = \sqrt{\lambda_{\mathrm{diff}}}$,
but that relation is a consequence of Gaussian structure and not a
general isometry. For non-abelian gauge theory no such correspondence is
available here.

The companion obstruction concerns the coarse-graining itself. A scale
map built from a gauge-covariant Markov kernel that preserves both
reflection positivity and locality is not available: any such kernel
either forces the holonomies to commute --- reducing the theory to an
abelian one --- or destroys reflection positivity, so that the
reconstructed inner product is indefinite. This is why the scale route
actually used in \S\ref{sec:continuum} is a \emph{deterministic}
pushforward rather than a Markov kernel, and why the two-spin
counterexample $\mathbb{E}[s_- s_+] = -1$ is recorded as a proposition
rather than as an anecdote.

\begin{scopenote}
Break~6 does not refute the curvature computations themselves. The Haar
Ricci constant and the Wilson Hessian bound are correct statements about
the configuration manifold. What is refuted is their use as a route to
$\Delta_{\mathrm{phys}}$, and the curvature results are consequently
recorded in this program as inputs to coercivity estimates rather than
as a mass-gap mechanism.
\end{scopenote}

\subsection{What survives}

\begin{enumerate}
  \item The exact strong-coupling spectral theory of
    Part~\ref{part:exact}, in its entirety. None of the six obstructions
    touches a fixed-order statement.
  \item The fixed-spacing infinite-volume Wilson band and its literal
    source frame (\S\ref{sec:wilson}), on the interval
    \eqref{eq:break5-u}, with the scope stated there.
  \item The moving-time spectral exclusion as an abstract interface: the
    implication from its hypotheses to a limiting gap is proved and
    sharp. Break~3 is a statement about whether its hypotheses hold for
    Yang--Mills, not about the implication.
\end{enumerate}

What does not survive is any inference from the above to a continuum
Yang--Mills mass gap. The program states that plainly, and
\S\ref{sec:obligations} is the list of what would have to change.
